\documentclass[11pt,a4paper]{article}
\usepackage[margin=19mm,headheight=15pt]{geometry}
\usepackage{fontspec}
\setmainfont{Noto Serif}
\usepackage{amsmath,amssymb,mathtools}
\usepackage{xcolor}
\usepackage{fancyhdr}
\usepackage{enumitem}
\usepackage{microtype}
\setlength{\parindent}{0pt}
\setlength{\parskip}{6pt}
\definecolor{navy}{HTML}{183153}
\definecolor{lightgray}{HTML}{F3F5F7}
\pagestyle{fancy}\fancyhf{}
\lhead{\small ME 5309 -- Mid-Semester Examination}
\rhead{\small ME 5309}
\cfoot{\thepage}
\begin{document}
\begin{center}
{\Large\bfseries ME 5309 -- Mid-Semester Examination}\\[3pt]
{\large Answer Booklet}\\[8pt]
\end{center}
\textbf{Instructions:} Attempt any three questions. Each question carries 10 marks. Show working where appropriate.
\section*{Question 1 \hfill [10 marks]}
\subsection*{Q1(a) \hfill [4 marks]}

Pixel pitch = 0.005 mm, sensor width = 1280x0.005 = 6.4 mm. how to calculate the belt field of view.

\subsection*{Q1(b) \hfill [3 marks]}

\[\begin{aligned}Sensitivity = 10 mV/N\end{aligned}\]
The ADC has 12 bits, so there are 4096 levels.

\subsection*{Q1(c) \hfill [3 marks]}

\[\begin{aligned}F = e^(AT) = e^(-0.2) = 0.818731\end{aligned}\]

\section*{Question 2 \hfill [10 marks]}
\subsection*{Q2(a) \hfill [2 marks]}

The state model is ẋ₁ = x₂, ẋ₂ = (τ-0.05x₂-2.4525 sin x₁)/0.02.
Equilibrium requires τ₀ = 2.4525 sin θ\_e, so τ₀ must be between -2.4525 and +2.4525 N m.

\subsection*{Q2(b) \hfill [2 marks]}

\[\begin{aligned}A = [[0,1], [-mg\ell sin(60^{\circ})/J, -b/J]] = [[0,1], [-106.2, -2.5]],\\ B = [[0],[50]]\end{aligned}\]

\subsection*{Q2(c) \hfill [4 marks]}

At rest, θ̇ = θ̈ = 0, so the torque balance is Kp(θd-θss) = mgℓ sin θss.
With Kp = 50, θd = π/2 and mgℓ = 2.4525, solve 50(π/2-θss) = 2.4525 sin θss.
This gives θss approximately 1.52181 rad = 87.193°, so the error is 90°-87.193° = 2.807°; it is not zero.\par\smallskip Rearranging the steady-state equation gives $50(\pi/2-\theta_{ss})=2.4525\sin\theta_{ss}$. I used radians in the equation. Solving numerically gives $\theta_{ss}\approx1.522\,$rad $\approx87.19^\circ$, so the gap from $90^\circ$ is about $2.81^\circ$.

\subsection*{Q2(d) \hfill [2 marks]}

With gravity compensation, the gravity term cancels the plant's mgℓ sin θ term.
At rest, Kp(θd-θss)=0, so θss=θd and the steady-state error is zero.
The linearised gravity stiffness term is also cancelled, so the closed-loop poles generally change; the stiffness contribution mgℓ cos θd is removed.\par\smallskip The compensation term cancels the nominal gravity torque in the plant model. At steady state the proportional term must then be zero, which requires $\theta_d-\theta_{ss}=0$. This does not mean the transient response is unchanged: removing the gravity derivative changes the linearized stiffness.

\section*{Question 4 \hfill [10 marks]}
\subsection*{Q4(a) \hfill [2 marks]}

The closed-loop poles are -2±2j.
I use Nff=c/b=6 for unity DC gain.

\subsection*{Q4(b) \hfill [2 marks]}

The nominal controller has unity DC gain, so θss=1 and the tracking error is zero even after the actuator gain changes.

\subsection*{Q4(c) \hfill [4 marks]}

The desired polynomial is (s+4)(s²+4s+8).
The closed-loop polynomial should be cubic because an integrator is added.

\subsection*{Q4(d) \hfill [2 marks]}

Integral action removes steady-state error, although this only holds for the nominal values of b and c.
Saturation can cause windup, which should be handled by anti-windup.

\end{document}